% GENERATED by scripts/manual_tables.py from the adapter registry. Do not edit:
% rerun the script. Hand-written text: blocks/detection/deconvolve/
\ubblock{detection.deconvolve}%
  {Deconvolve (OASIS)}% what the app calls it
  {Deconvolve: sudden jumps, each followed by the same fade (OASIS)}% registry display name
  {detect}% insert-modal tab
  {Signal}% input
  {Scores}% output
  {Adapters/detection\_deconvolve.py}% code
  {detection/deconvolve}% hand-written folder under blocks/
\ubblockdescription{detection.deconvolve}{Explains the trace as sudden jumps, each followed by the same exponential fade, and scores every sample by the size of its jump in noise SDs (almost all zero). A drop followed by a recovery is the shape upside down; a second drop that starts before the first has recovered is a jump of its own, which a threshold would merge. The baseline is one constant per span unless baseline\_window\_s is set: put a long Baseline removal (or BEADS) first and/or set baseline\_window\_s, and Impulse events after it. The recovery half-time is the first knob to set: read it off a few of the recording's own troughs (the time a trough takes to climb half-way back); 0 estimates it, and a wrong one fragments every trough into many jumps, which the summary flags. Then min\_jump\_sigma. Events that are not a sudden drop followed by an exponential recovery -- a sharkfin's plateau and slow climb -- are split into many small events; check the residual SD: above \textasciitilde{}1.5 the event count is not a trough count. The noise SD (the score's unit) is by default the background's own spread about its baseline (noise\_scale = level), which slow, coloured noise does not fool; noise\_scale = diff reads it from sample-to-sample differences, as the paper does, and lets coloured noise's wander pass as jumps. When events fill the trace with no quiet stretch between them, the level SD measures the events: compare sigma\_level with sigma in the readout. Friedrich, Zhou \& Paninski 2017 (OASIS), written from the paper.}
\ubblockflags{detection.deconvolve}{%
  \ubflag{estimate}{}%
}
\ubblocksideinputs{detection.deconvolve}{%
}
\ubblockparams{detection.deconvolve}{%
  \ubparam{polarity}{\texttt{drop}}{one of \texttt{drop}, \texttt{spike}}{drop: a sudden fall then a recovery (the signal is negated); spike: a sudden rise then a decay}%
  \ubparam{recovery\_\allowbreak{}half\_\allowbreak{}s}{0.0}{number, at least 0}{Seconds for the fade to recover half-way -- the FIRST knob: read it off a few of the recording's own troughs. 0 estimates it from the span (reported in the readout, with a warning when it cannot be trusted)}%
  \ubparam{rise\_\allowbreak{}s}{0.0}{number, at least 0}{Rise time constant in seconds for a jump that is not instant (makes the model AR(2)); 0 = an instant jump (AR(1)). Must be shorter than the recovery}%
  \ubparam{min\_\allowbreak{}jump\_\allowbreak{}sigma}{3.0}{number from 0 to 1000}{The smallest jump the fit may use, in noise SDs (the paper's s\_min; the SD noise\_scale names): raise it to stop noise being explained as tiny jumps}%
  \ubparam{baseline\_\allowbreak{}window\_\allowbreak{}s}{0.0}{number, at least 0}{0 = one constant baseline per span (the paper). > 0 = a baseline that follows slow drift, estimated over windows of this many seconds; make it several recoveries long (e.g. 15x the recovery half-time) so it follows the drift, not the events}%
  \ubparam{noise\_\allowbreak{}scale}{\texttt{level}}{one of \texttt{level}, \texttt{diff}}{The noise SD the score and min\_jump\_sigma are measured in. level: the SD of the background about its baseline (troughs clipped away) -- right for slow, coloured noise, where each sample sits close to the last. diff: from sample-to-sample differences (the earlier behaviour) -- reads coloured noise far too small, so its wander passes as jumps}%
}
